% GENERATED FILE. Edit canonical lessons, then run npm run generate:books.
% canonical-source-hash: fnv1a64:7c4d51807d10f05f
% canonical-lessons: GE-C65-warten, GE-C65-bezahlen, GE-C65-kochen, GE-C65-lachen, GE-C65-weinen

\chapter{To wait, To pay, To cook, To laugh, To cry}
\label{ch:ge-to-wait-to-pay-to-cook-to-laugh-to-cry}
\hlchaptermodality{65}

\begin{chapteropening}
\textbf{By the end of this chapter:} \emph{I can say to wait, to pay, to cook, to laugh and to cry.}

Complete the last lesson of chapter 65: I can say to wait, to pay, to cook, to laugh and to cry.
\end{chapteropening}

\section[warten]{warten — to wait}

\label{lesson:GE-C65-warten}



\begin{quote}
Before the new one: say the German for to find, then the German for to look for.
\end{quote}

\subsection*{You'll want to know: warten}
\textbf{warten} — \textquotedblleft{}to wait\textquotedblright{}.

\begin{cousinweb}
\textbf{warten} $\leftarrow$ Old High German \emph{wartēn}, \textquotedblleft{}to watch, to look after\textquotedblright{} $\leftarrow$ Proto-Germanic \emph{*wardāną}, \textquotedblleft{}to guard\textquotedblright{} — inherited, a cousin of English \emph{ward}. (EtymDB 2.1) Waiting began as keeping watch.
\end{cousinweb}

\begin{grammarlens}[title={What you've built}]
One more word for this chapter.
\end{grammarlens}

\subsection*{Your turn}
\begin{itemize}
\raggedright
  \item \emph{Say it:} \emph{warten}
  \item \emph{Say it:} \emph{warten}, once more
  \item \emph{Say it:} say \emph{warten}
  \item \emph{Recall:} say the German for to find, then the German for to look for, then say \emph{warten} again
\end{itemize}

\begin{tcolorbox}[breakable,colback=teal!4,colframe=teal!35!black,title={Before you move on}]
What does warten mean? (\textquotedblleft{}To wait\textquotedblright{}.) Say it once more.
\end{tcolorbox}
\section[bezahlen]{bezahlen — to pay}

\label{lesson:GE-C65-bezahlen}



\begin{quote}
Before the new one: say the German for to look for, then the German for to wait.
\end{quote}

\subsection*{You'll want to know: bezahlen}
\textbf{bezahlen} — \textquotedblleft{}to pay\textquotedblright{}.

\begin{cousinweb}
\textbf{bezahlen} is \emph{be-} + the verb \textquotedblleft{}to pay\textquotedblright{}, which is built on the word for \textquotedblleft{}number\textquotedblright{}, Old High German \emph{zala} $\leftarrow$ Proto-Germanic \emph{*talō} (EtymDB 2.1). English \emph{tale} and \emph{tell} share the root: paying was counting out.
\end{cousinweb}

\begin{grammarlens}[title={What you've built}]
One more word for this chapter.
\end{grammarlens}

\subsection*{Your turn}
\begin{itemize}
\raggedright
  \item \emph{Say it:} \emph{bezahlen}
  \item \emph{Say it:} \emph{bezahlen}, once more
  \item \emph{Say it:} say \emph{bezahlen}
  \item \emph{Recall:} say the German for to look for, then the German for to wait, then say \emph{bezahlen} again
\end{itemize}

\begin{tcolorbox}[breakable,colback=teal!4,colframe=teal!35!black,title={Before you move on}]
What does bezahlen mean? (\textquotedblleft{}To pay\textquotedblright{}.) Say it once more.
\end{tcolorbox}
\section[kochen]{kochen — to cook}

\label{lesson:GE-C65-kochen}



\begin{quote}
Before the new one: say the German for to wait, then the German for to pay.
\end{quote}

\subsection*{You'll want to know: kochen}
\textbf{kochen} — \textquotedblleft{}to cook\textquotedblright{}.

\begin{cousinweb}
\textbf{kochen} $\leftarrow$ Old High German \emph{kohhōn} $\leftarrow$ Latin \emph{coquere}, \textquotedblleft{}to cook\textquotedblright{} (EtymDB 2.1) — borrowed from Roman kitchens, as English \emph{cook} was.
\end{cousinweb}

\begin{grammarlens}[title={What you've built}]
One more word for this chapter.
\end{grammarlens}

\subsection*{Your turn}
\begin{itemize}
\raggedright
  \item \emph{Say it:} \emph{kochen}
  \item \emph{Say it:} \emph{kochen}, once more
  \item \emph{Say it:} say \emph{kochen}
  \item \emph{Recall:} say the German for to wait, then the German for to pay, then say \emph{kochen} again
\end{itemize}

\begin{tcolorbox}[breakable,colback=teal!4,colframe=teal!35!black,title={Before you move on}]
What does kochen mean? (\textquotedblleft{}To cook\textquotedblright{}.) Say it once more.
\end{tcolorbox}
\section[lachen]{lachen — to laugh}

\label{lesson:GE-C65-lachen}



\begin{quote}
Before the new one: say the German for to pay, then the German for to cook.
\end{quote}

\subsection*{You'll want to know: lachen}
\textbf{lachen} — \textquotedblleft{}to laugh\textquotedblright{}.

\begin{cousinweb}
\textbf{lachen} $\leftarrow$ Old High German \emph{lahhan} $\leftarrow$ Proto-Germanic \emph{*hlahjaną} — inherited, the same verb as English \emph{laugh}. (EtymDB 2.1)
\end{cousinweb}

\begin{grammarlens}[title={What you've built}]
One more word for this chapter.
\end{grammarlens}

\subsection*{Your turn}
\begin{itemize}
\raggedright
  \item \emph{Say it:} \emph{lachen}
  \item \emph{Say it:} \emph{lachen}, once more
  \item \emph{Say it:} say \emph{lachen}
  \item \emph{Recall:} say the German for to pay, then the German for to cook, then say \emph{lachen} again
\end{itemize}

\begin{tcolorbox}[breakable,colback=teal!4,colframe=teal!35!black,title={Before you move on}]
What does lachen mean? (\textquotedblleft{}To laugh\textquotedblright{}.) Say it once more.
\end{tcolorbox}
\section[weinen]{weinen — to cry}

\label{lesson:GE-C65-weinen}



\begin{quote}
Before the new one: say the German for to cook, then the German for to laugh.
\end{quote}

\subsection*{You'll want to know: weinen}
\textbf{weinen} — \textquotedblleft{}to cry\textquotedblright{}.

\begin{cousinweb}
\textbf{weinen} $\leftarrow$ Old High German \emph{weinōn} $\leftarrow$ Proto-Germanic \emph{*wainōną}, built on the cry \emph{wai}, \textquotedblleft{}woe\textquotedblright{} — inherited, and unrelated to \emph{Wein}. (EtymDB 2.1)
\end{cousinweb}

\begin{grammarlens}[title={What you've built}]
One more word for this chapter.
\end{grammarlens}

\subsection*{Your turn}
\begin{itemize}
\raggedright
  \item \emph{Say it:} \emph{weinen}
  \item \emph{Say it:} \emph{weinen}, once more
  \item \emph{Say it:} say \emph{weinen}
  \item \emph{Recall:} say the German for to cook, then the German for to laugh, then say \emph{weinen} again
\end{itemize}

\begin{tcolorbox}[breakable,colback=teal!4,colframe=teal!35!black,title={Before you move on}]
What does weinen mean? (\textquotedblleft{}To cry\textquotedblright{}.) Say it once more.
\end{tcolorbox}
